\chapter{Volatility as a Traded Quantity: First Contact}\label{ch:m1:volatility-first-contact}

The three-month 100 call is offered at 4.06. Nobody on an options desk says
so. They say the call is ``twenty vol''; that the 90 put is ``twenty-four'';
that the stock's ``skew is six points''; that ``the forward is 100.20, so the
market has an 80-cent dividend in''. Prices in dollars change every time the
share ticks; these numbers do not, and so they are what the desk quotes,
compares, remembers and trades. This chapter learns the language. It needs
one formula from outside (the Black--Scholes price, derived in One Quant
Book 5 and used here as a black box that turns a volatility into a price)
and one identity that needs no model at all.

\section{Put--call parity, with dividends and borrow}

\begin{proposition}[Put--call parity]\label{prop:m1:volatility-first-contact:parity}
For European options with the same strike $K$ and expiry $T$ on the same
underlying,
\[
C - P \;=\; \mathrm{e}^{-rT}\,(F - K),
\]
where $F$ is the forward price of the underlying for $T$: the \emph{put--call
parity}\index{put--call parity}. For a share, $F\,\mathrm{e}^{-rT} = S - D$,
with $D$ the present value of the dividends paid before $T$, to which a
short seller adds the borrow fee.
\end{proposition}
\begin{proof}
Long a call and short a put at $K$ is an obligation to buy at $K$ at $T$: a
forward contract struck at $K$, worth $\mathrm{e}^{-rT}(F-K)$ today.
\end{proof}

Parity holds whatever the volatility, the distribution or the model. It is
enforced by two trades that a \omterm{def:m1:what-a-trading-firm-does:market-maker}{market maker} can do mechanically: the
\emph{conversion} (buy the share, buy the put, sell the call) and the
\emph{reversal} (the opposite, which needs a stock borrow). Read backwards,
it is a measuring instrument:

\begin{method}[Reading the forward off the chain]\label{met:m1:volatility-first-contact:forward}
\begin{enumerate}
\item At each strike compute $F_K = K + \mathrm{e}^{rT}(C_K - P_K)$ from
mid-prices. The values should agree to within the quotes' width; average the
ones near the money.
\item The implied carry is $D^{\text{imp}} = S - \mathrm{e}^{-rT}F$: expected
dividends plus the borrow fee that the marginal arbitrageur pays.
\item If $D^{\text{imp}}$ exceeds any plausible dividend, the share is hard
to borrow (\cref{ch:m1:stock-loan-and-short-selling}) and the options are
telling you the fee.
\end{enumerate}
\end{method}

For American options parity is an inequality, since both sides can be
exercised early; near the money and away from dividends the European relation
is a close approximation, and it fails in exactly the cases of
\cref{ch:m1:option-contracts-and-exchanges} in which early exercise is
rational.

\section{Implied volatility}

\begin{definition}[Implied volatility]\label{def:m1:volatility-first-contact:iv}
The \emph{implied volatility}\index{implied volatility} of an option is the
value of the volatility parameter which, put into the Black--Scholes
formula with the option's strike, expiry, forward and interest rate, gives
the option's market price. It is a change of units, from dollars to annualised
standard deviation of returns, and asserts nothing about the model's truth.
\end{definition}

\begin{definition}[At-the-money]\label{def:m1:volatility-first-contact:atm}
An option is \emph{at-the-money}\index{at-the-money} when its strike equals
the forward (or, loosely, the spot). Its price is, to a good approximation,
proportional to volatility:
$C_{\text{ATM}} \approx 0.4\,F\,\mathrm{e}^{-rT}\sigma\sqrt{T}$.
\end{definition}

An option's price is increasing in volatility, from its discounted intrinsic
value at $\sigma = 0$ to the discounted forward (call) or strike (put) as
$\sigma \to \infty$. Between those bounds there is exactly one \omterm{def:m1:volatility-first-contact:iv}{implied
volatility}, and bisection finds it; outside them there is none, and a
solver must say so: a price below intrinsic value is a stale quote or a
wrong forward, not a negative volatility.

\begin{omfigure}
\begin{tikzpicture}
\begin{axis}[width=10.5cm,height=4.8cm,
  xlabel={volatility (\% a year)},ylabel={option price},
  tick label style={font=\small},label style={font=\small},
  xmin=0,xmax=60,ymin=0,ymax=13,grid=major,grid style={gray!25},
  legend columns=2,legend style={font=\footnotesize,at={(0.5,-0.32)},anchor=north,draw=none,fill=none,/tikz/every even column/.append style={column sep=8pt}},
  table/col sep=comma]
\addplot[omDef,very thick] table[x=vol_pct,y=call_100]{figdata/markets-1/25-volatility-first-contact/price_vs_vol.csv};
\addplot[omThm,very thick,dashed] table[x=vol_pct,y=put_90]{figdata/markets-1/25-volatility-first-contact/price_vs_vol.csv};
\draw[thin] (axis cs:0,4.06) -- (axis cs:20.08,4.06) -- (axis cs:20.08,0);
\draw[thin] (axis cs:0,1.12) -- (axis cs:23.82,1.12) -- (axis cs:23.82,0);
\legend{100 call,90 put}
\end{axis}
\end{tikzpicture}

\omcaption{Price against volatility for two three-month options on a share at
100 with a forward of 100.20. Inverting the curve at the market price gives
the \omterm{def:m1:volatility-first-contact:iv}{implied volatility}: 4.06 is 20.1 for the call, 1.12 is 23.8 for the put.
Near the money the curve is a straight line; far from it, it is flat at low
volatility, which is where solvers lose precision. Data: the chapter's
build.}\label{fig:m1:volatility-first-contact:inversion}
\end{omfigure}

\begin{definition}[Realised volatility and variance]\label{def:m1:volatility-first-contact:realised}
The \emph{realised volatility}\index{realised volatility} of an asset over a
period is the annualised standard deviation of its returns over that
period, conventionally computed from daily log returns without subtracting
the mean. \emph{Variance}\index{variance} is its square. \omterm{def:m1:volatility-first-contact:iv}{Implied volatility} is
a price set today; realised volatility is a statistic known afterwards; the
difference between them is what a hedged option position earns or loses.
\end{definition}

\section{Skew and term structure}

\begin{definition}[Volatility skew and term structure]\label{def:m1:volatility-first-contact:skew}
The \emph{volatility skew}\index{volatility skew} (or smile) is the
dependence of \omterm{def:m1:volatility-first-contact:iv}{implied volatility} on strike for one expiry; the \emph{term
structure of volatility}\index{term structure of volatility} is its
dependence on expiry for a given moneyness. Together they form the
volatility surface.
\end{definition}

If the model behind the change of units were true, every strike would have
the same \omterm{def:m1:volatility-first-contact:iv}{implied volatility}. For equity indices since 1987 the \omterm{def:m1:volatility-first-contact:iv}{implied
volatility} of low strikes has been well above that of high strikes: crashes
are more frequent and more feared than the model allows, and investors pay
for protection. The skew is quoted as a difference in volatility points
between two reference strikes or deltas, and traded as such
(\cref{fig:m1:volatility-first-contact:smile}).

\begin{omfigure}
\begin{tikzpicture}
\begin{axis}[width=10.5cm,height=4.8cm,
  xlabel={strike},ylabel={implied volatility (\%)},
  tick label style={font=\small},label style={font=\small},
  xmin=78,xmax=122,ymin=14,ymax=33,grid=major,grid style={gray!25},
  nodes near coords,nodes near coords style={font=\scriptsize,above=2pt,/pgf/number format/.cd,fixed,fixed zerofill,precision=1},
  point meta=explicit,table/col sep=comma]
\addplot[omDef,very thick,mark=*] table[x=strike,y=implied_vol_pct,meta=implied_vol_pct]{figdata/markets-1/25-volatility-first-contact/chain.csv};
\end{axis}
\end{tikzpicture}

\omcaption{The smile recovered from the screen prices of the chapter's
synthetic chain, out-of-the-money option at each strike, with the forward
read from parity. The 90--110 skew is 6.2 volatility points. Data: the
tutorial.}\label{fig:m1:volatility-first-contact:smile}
\end{omfigure}

\section{The volatility index}

\begin{definition}[Volatility index]\label{def:m1:volatility-first-contact:vix}
A \emph{volatility index}\index{volatility index} is a published number that
summarises the \omterm{def:m1:volatility-first-contact:iv}{implied volatility} of an index's options for a fixed horizon,
computed from option prices by a formula that uses no pricing model:
\[
\sigma^2 \;=\; \frac{2}{T}\sum_i \frac{\Delta K_i}{K_i^2}\,\mathrm{e}^{RT}\,Q(K_i)
\;-\; \frac{1}{T}\Bigl(\frac{F}{K_0}-1\Bigr)^2 ,
\]
with $Q(K_i)$ the mid-quote of the out-of-the-money option at strike $K_i$,
$K_0$ the first strike at or below the forward $F$, and the index equal to
$100\,\sigma$.
\end{definition}

\begin{dated}{2026-09}{The index and its futures}\label{dat:m1:volatility-first-contact:vix}
\emph{The index.} The best-known \omterm{def:m1:volatility-first-contact:vix}{volatility index} uses S\&P 500 options of
two expiries bracketing 30 days and interpolates their \omterm{def:m1:volatility-first-contact:realised}{variances} to a constant
30-day horizon. Only options with a non-zero bid enter, at mid-quote; going
away from the money, the strip stops once two consecutive strikes have no
bid. $K_0$ is the strike at or immediately below the forward, itself found
from the strike where call and put prices are closest. \emph{The futures.}
Futures on the index have a multiplier of \$1\,000 and a tick of 0.05 point,
\$50. They settle on the Wednesday thirty days before the third Friday of the
following month, against a special opening quotation of the index computed
from the opening prices of the S\&P 500 options that morning.
\end{dated}

The sum weights each option by $1/K^2$: low strikes count for more, so the
index rises with the skew as well as with the general level of \omterm{def:m1:volatility-first-contact:iv}{implied
volatility}. On the synthetic chain of this chapter the \omterm{def:m1:volatility-first-contact:atm}{at-the-money}
volatility is 20.1 and the index-style number 23.0
(\cref{fig:m1:volatility-first-contact:weights}). The formula is the price of
the portfolio of options that replicates the \emph{\omterm{def:m1:volatility-first-contact:realised}{variance}} the index will
realise, a result proved in One Quant Book 5; here it matters that the
index is a price of a traded portfolio, not an opinion.

\begin{omfigure}
\begin{tikzpicture}
\begin{axis}[width=10cm,height=4.4cm,ybar,bar width=6pt,
  xlabel={strike (forward at 100)},ylabel={weight (\%)},
  tick label style={font=\small},label style={font=\small},
  xmin=55,xmax=145,ymin=0,ymax=3.4,ymajorgrids,grid style={gray!25},table/col sep=comma]
\addplot[fill=omDef!60,draw=omDef] table[x=strike,y=weight]{figdata/markets-1/25-volatility-first-contact/weights.csv};
\end{axis}
\end{tikzpicture}

\omcaption{The weight $\Delta K/K^2$ given to one dollar of \omterm{def:m1:option-contracts-and-exchanges:option}{option premium} at
each strike, normalised. A dollar of premium in the 60 put counts five times
as much as a dollar in the 140 call. Data: the
definition.}\label{fig:m1:volatility-first-contact:weights}
\end{omfigure}

\section{Futures on volatility}

\begin{definition}[Volatility-index future]\label{def:m1:volatility-first-contact:vixfut}
A \emph{VIX future}\index{VIX future} is a \omterm{def:m1:basis-roll-and-delivery:delivery}{cash-settled} future on the value
of the \omterm{def:m1:volatility-first-contact:vix}{volatility index} at its expiry. Since the index cannot be bought and
held, the future is not tied to today's index by a \omterm{def:m1:basis-roll-and-delivery:carry}{cost of carry}: it is the
market's price for where the index will be.
\end{definition}

Two facts organise this market. Volatility mean-reverts, so the futures
curve usually slopes up when the index is low (\omterm{def:m1:basis-roll-and-delivery:contango}{contango}) and down when it has
spiked (\omterm{def:m1:basis-roll-and-delivery:contango}{backwardation}). And sellers of volatility are paid a premium on
average: in calm periods a position short the front future earns the roll
down the curve month after month, like the \omterm{def:m1:basis-roll-and-delivery:contango}{contango} of
\cref{fig:m1:basis-roll-and-delivery:rolled} seen from the other side.

\begin{omfigure}
\begin{tikzpicture}
\begin{axis}[width=10.5cm,height=4.8cm,
  xlabel={months to expiry (0 = the index)},ylabel={volatility points},
  tick label style={font=\small},label style={font=\small},
  xmin=0,xmax=7,ymin=10,ymax=40,grid=major,grid style={gray!25},
  legend columns=2,legend style={font=\footnotesize,at={(0.5,-0.32)},anchor=north,draw=none,fill=none,/tikz/every even column/.append style={column sep=8pt}},
  table/col sep=comma]
\addplot[omDef,very thick,mark=*] table[x=month,y=calm]{figdata/markets-1/25-volatility-first-contact/term.csv};
\addplot[omThm,very thick,dashed,mark=square*] table[x=month,y=stress]{figdata/markets-1/25-volatility-first-contact/term.csv};
\legend{a calm market: contango,after a spike: backwardation}
\end{axis}
\end{tikzpicture}

\omcaption{Two stylised curves of futures on a \omterm{def:m1:volatility-first-contact:vix}{volatility index}. In the calm
curve a short position in the first future earns 1.6 points a month if
nothing happens; in the other, the market expects the index to fall by a
third within three months. Illustrative
numbers.}\label{fig:m1:volatility-first-contact:term}
\end{omfigure}

\begin{example}[5 February 2018]\label{ex:m1:volatility-first-contact:feb2018}
On Monday 5 February 2018 the S\&P 500 fell 4\% and the \omterm{def:m1:volatility-first-contact:vix}{volatility index} rose
20 points in the day. Exchange-traded products giving
leveraged or inverse exposure to the front futures, with about \$4 billion of
assets at the end of 2017, had to rebalance at the end of the day, and, by
the arithmetic of \cref{prop:m1:exchange-traded-funds:rebalance}, both kinds
had to \emph{buy} futures after a rise. A central-bank analysis describes the
resulting loop: their buying pushed the futures higher, which increased the
amount they had to buy. The largest inverse product lost nearly all of its
value in that session and was then terminated by its issuer.
\end{example}

\section{Tutorial: from a screen of prices to a forward, a dividend and a smile}\label{tut:m1:volatility-first-contact:read}

\begin{tutorial}
\textbf{Goal.} Price with the black-box formula, invert it robustly, read the
forward and the implied dividend from parity, recover the smile, and compute
an index-style volatility. \textbf{End state:} the four data figures of this
chapter.
\begin{tutsteps}
\item \textbf{The solver.} Bounds first, then bisection.
\omcode{code/firm/parity/firm_parity.py}{24}{37}{Implied volatility, or an error when no volatility can explain the price.}\label{lst:m1:volatility-first-contact:solver}
\item \textbf{Parity as an instrument.}
\omcode{code/firm/parity/firm_parity.py}{40}{48}{The forward from one strike, and the carry it implies.}\label{lst:m1:volatility-first-contact:parity}
\item \textbf{Reading a chain} rounded to the cent, as a screen shows it.
\omcode{code/markets-1/25-volatility-first-contact/python/vol_first.py}{33}{42}{Forward, dividend and smile from prices alone.}\label{lst:m1:volatility-first-contact:read}
\item \textbf{The strip.}
\omcode{code/firm/parity/firm_parity.py}{65}{80}{Model-free implied variance from out-of-the-money quotes.}\label{lst:m1:volatility-first-contact:strip}
\end{tutsteps}
\textbf{What to change next.} Truncate the strip at 80 and 120 and see how
much of the index-style volatility is lost; then widen the strike spacing
from 1 to 5 and measure the discretisation error. Both are real features of
published indices.
\end{tutorial}

\section{Build: parity checker and implied-volatility solver}\label{bld:m1:volatility-first-contact:parity}

\begin{build}
\bfield{Purpose} The first two functions every options system of the
miniature firm calls: what volatility is this price, and is this chain
consistent with one forward?
\bfield{Interface} \texttt{price(forward, strike, years, rate, vol, right)};
\texttt{implied\_vol(premium, forward, strike, years, rate, right)};
\texttt{implied\_forward(call, put, strike, years, rate)};
\texttt{implied\_dividends\_pv(spot, forward, years, rate)};
\texttt{conversion\_edge(\ldots)} and \texttt{reversal\_edge(\ldots)} from bid
and ask prices; \texttt{variance\_strip(forward, years, rate, quotes)}.
\bfield{Rules} The solver raises on a premium outside the no-arbitrage
bounds and never returns a negative or a clipped volatility silently.
Conversions and reversals are evaluated at the prices one would actually
trade: call bid and put ask for the one, the reverse for the other, with the
borrow cost from \cref{bld:m1:stock-loan-and-short-selling:borrow} on the
short side. Series and expiries come from the chain container
(\cref{bld:m1:option-contracts-and-exchanges:chain}).
\bfield{Acceptance tests} \texttt{code/firm/parity/tests/}: parity of the
pricer; round trips of the solver across strikes and volatilities, and its
refusals; a dividend recovered from a forward; conversions and reversals
unprofitable inside the quotes and profitable on a planted error; the strip
recovering a flat volatility to two hundredths of a point.
\bfield{Stretch} American options: an early-exercise premium by a binomial
tree, and the inequality form of parity.
\end{build}

\begin{omsources}
\item Cboe Global Indices, \emph{Volatility Index Methodology: Cboe
Volatility Index} (formula, selection of options, interpolation).
\item Cboe Futures Exchange, \emph{VIX Futures: contract specifications}.
\item Bank for International Settlements, ``The equity market turbulence of 5
February: the role of exchange-traded volatility products'', \emph{BIS
Quarterly Review}, March 2018.
\item F.~Black and M.~Scholes, ``The pricing of options and corporate
liabilities'', \emph{Journal of Political Economy} 81 (1973); F.~Black,
``The pricing of commodity contracts'', \emph{Journal of Financial Economics}
3 (1976).
\end{omsources}

\section{Exercises}

\begin{exercise}[$\star$]\label{exo:m1:volatility-first-contact:1}
With $r = 4\%$ and $T = 0.25$, the 100 call is at 4.06 and the 100 put at 3.87.
Give the implied forward. With the share at 100, give the implied carry.
\end{exercise}

\begin{exercise}[$\star$]\label{exo:m1:volatility-first-contact:2}
Use the \omterm{def:m1:volatility-first-contact:atm}{at-the-money} approximation to estimate the \omterm{def:m1:volatility-first-contact:iv}{implied volatility} of the
100 call of the previous exercise, and compare with the solver's 20.1.
\end{exercise}

\begin{exercise}[$\star$]\label{exo:m1:volatility-first-contact:3}
A volatility-index future is bought at 16.70 and sold at 19.25. Give the
P\&L of 40 contracts. How many ticks is that?
\end{exercise}

\begin{exercise}[$\star\star$]\label{exo:m1:volatility-first-contact:4}
A three-month 90 call on a share at 100 (forward 100.20, $r = 4\%$) is quoted
at 10.05. Show that it has no \omterm{def:m1:volatility-first-contact:iv}{implied volatility}. What should a system do
with this quote?
\end{exercise}

\begin{exercise}[$\star\star$]\label{exo:m1:volatility-first-contact:5}
From a chain you read a forward of 98.10 for three months with the share at
100 and $r = 4\%$. The company pays no dividend. Give the implied borrow fee
as an annual rate and interpret it.
\end{exercise}

\begin{exercise}[$\star\star$]\label{exo:m1:volatility-first-contact:6}
An index's daily log returns over five days are $+0.4\%$, $-1.1\%$, $+0.7\%$,
$-2.0\%$, $+1.5\%$. Give the \omterm{def:m1:volatility-first-contact:realised}{realised volatility}, annualised with 252 days and
without subtracting the mean. If \omterm{def:m1:volatility-first-contact:iv}{implied volatility} was 16, who was right?
\end{exercise}

\begin{exercise}[$\star\star\star$]\label{exo:m1:volatility-first-contact:7}
\emph{Coding.} With \texttt{read\_chain} on the chapter's nine-strike chain
report the average implied forward, the range of the nine strike-by-strike
forwards, the implied dividend and the 90--110 skew. Then round prices to
five cents instead of one and report the same: which numbers survive?
\end{exercise}

\begin{exercise}[$\star\star\star$]\label{exo:m1:volatility-first-contact:8}
\emph{Find the flaw.} ``The \omterm{def:m1:volatility-first-contact:vix}{volatility index} is at 14 and its long-run average
is 19. It is cheap: buy the front future at 15.6 and wait for mean
reversion.'' Use \cref{fig:m1:volatility-first-contact:term}.
\end{exercise}

\section{Problem: Reading a Chain}

\begin{problem}[{Weekend problem --- everything the screen knows}]\label{pb:m1:volatility-first-contact:1}
A share trades at 100.00. Three-month European options ($T = 0.25$, $r = 4\%$)
show these mid-prices:

\begin{center}\small
\begin{tabular}{lrrrrr}
\toprule
Strike & 90 & 95 & 100 & 105 & 110 \\
\midrule
Call & 11.22 & 7.26 & 4.06 & 1.88 & 0.68 \\
Put & 1.12 & 2.11 & 3.87 & 6.63 & 10.38 \\
\bottomrule
\end{tabular}
\end{center}

\medskip\noindent\textbf{Part I --- The forward.}
\begin{enumerate}
\item Compute the implied forward at each of the five strikes.
\item Why do they differ slightly? Which would you trust least?
\item Give their average.
\item Give the implied carry, in dollars.
\item The company is expected to pay 0.80 in five weeks. Is the share hard
to borrow?
\end{enumerate}

\medskip\noindent\textbf{Part II --- The smile.}
\begin{enumerate}[resume]
\item Which option do you use at each strike to compute \omterm{def:m1:volatility-first-contact:iv}{implied volatility},
and why?
\item The solver gives 23.8, 21.8, 20.1, 18.7 and 17.6. Give the 90--110 skew
and the slope in volatility points per 10\% of strike.
\item Check the \omterm{def:m1:volatility-first-contact:atm}{at-the-money} figure with the approximation of
\cref{def:m1:volatility-first-contact:atm}.
\item A colleague computes the \omterm{def:m1:volatility-first-contact:iv}{implied volatility} of the 90 \emph{call} with
the spot in place of the forward and finds 25.1 instead of 23.8. Explain.
\end{enumerate}

\medskip\noindent\textbf{Part III --- Trading it.}
\begin{enumerate}[resume]
\item The 100 call is bid 4.03 and offered 4.09; the 100 put 3.84 at 3.90; the
share 99.99 at 100.01. With the 0.80 dividend worth 0.797 today, evaluate the
conversion.
\item Evaluate the reversal, with a borrow cost of 0.05.
\item What do the two results say about the market?
\item A quote appears: 100 call bid at 4.60. What do you do, and what do you
check first?
\item The next day the share is at 97 and the 100 call at 2.74. Has the call
become ``cheaper''? What would you compute to answer?
\end{enumerate}

\medskip\noindent\textbf{Part IV --- Judgement.}
\begin{enumerate}[resume]
\item Why do desks quote volatility and not price?
\item What does a skew of six points say about the market's view of the
share?
\item On this chain an index-style volatility would be 23.0 against 20.1 at
the money. Why is it higher?
\item Why can \omterm{def:m1:volatility-first-contact:iv}{implied volatility} be above \omterm{def:m1:volatility-first-contact:realised}{realised volatility} on average
without anyone being irrational?
\item State the \emph{named result}: the implied dividend and the 90--110
skew.
\item In one sentence: what is \omterm{def:m1:volatility-first-contact:iv}{implied volatility}?
\end{enumerate}
\end{problem}

\section{Interview questions}

\begin{interviewq}[$\star$ \iqroles{trader, researcher}]\label{iq:m1:volatility-first-contact:1}
State \omterm{prop:m1:volatility-first-contact:parity}{put--call parity} and say what it assumes.
\end{interviewq}

\begin{interviewq}[$\star$ \iqroles{trader, researcher}]\label{iq:m1:volatility-first-contact:2}
What is \omterm{def:m1:volatility-first-contact:iv}{implied volatility}, and why is it not a forecast?
\end{interviewq}

\begin{interviewq}[$\star\star$ \iqroles{developer, researcher}]\label{iq:m1:volatility-first-contact:3}
Write an implied-volatility solver. What can go wrong?
\end{interviewq}

\begin{interviewq}[$\star\star$ \iqroles{trader, researcher}]\label{iq:m1:volatility-first-contact:4}
Calls look cheap relative to puts at the same strike. What are the possible
explanations before you call it an arbitrage?
\end{interviewq}

\begin{interviewq}[$\star\star$ \iqroles{researcher, trader}]\label{iq:m1:volatility-first-contact:5}
How is a \omterm{def:m1:volatility-first-contact:vix}{volatility index} computed, and why can you not buy it?
\end{interviewq}

\begin{interviewq}[$\star\star\star$ \iqroles{trader, researcher}]\label{iq:m1:volatility-first-contact:6}
Explain what happened to short-volatility products in February 2018, and
what a risk manager should have seen beforehand.
\end{interviewq}
